\documentclass[11pt,letterpaper]{article}
\usepackage[margin=1in]{geometry}
\usepackage{amsmath,amssymb}
\usepackage{booktabs}
\usepackage{hyperref}
\usepackage{xcolor}
\usepackage{listings}
\usepackage{graphicx}
\usepackage{array}
\usepackage{longtable}
\usepackage{fancyhdr}
\usepackage{titlesec}
\newcommand{\doi}[1]{\href{https://doi.org/#1}{doi:#1}}

\hypersetup{
  colorlinks=true,
  linkcolor=blue!70!black,
  urlcolor=blue!70!black,
  citecolor=green!50!black
}

\lstset{
  basicstyle=\ttfamily\small,
  backgroundcolor=\color{gray!10},
  frame=single,
  breaklines=true,
  keywordstyle=\color{blue},
  commentstyle=\color{green!60!black},
  stringstyle=\color{red!70!black},
}

\pagestyle{fancy}
\fancyhf{}
\rhead{ERF Instrument Simulators}
\lhead{Design Document}
\cfoot{\thepage}

\title{\textbf{Instrument Simulator Design Document}\\
\large ERF Atmospheric Model\\
\normalsize Ceilometer $\cdot$ Doppler Lidar $\cdot$ Microwave Radiometer}
\author{ERF Development Team}
\date{\today}

\begin{document}

\maketitle
\tableofcontents
\newpage

% ============================================================
\section{Overview}
% ============================================================

This document describes the physics, algorithms, and software design of three ground-based
remote sensing instrument simulators implemented in the ERF atmospheric model.
The simulators provide forward operators that translate model state variables into
synthetic instrument observables, enabling direct comparison with field campaign measurements.

All three simulators are implemented as header-only C++ structs in
\texttt{Source/IO/ERF\_InstrumentSimulators.H}.
They are activated via \texttt{ParmParse} parameters in the input file and
written on a user-specified timestep interval.

\subsection{Instruments}

\begin{center}
\begin{tabular}{llll}
\toprule
\textbf{Instrument} & \textbf{Make/Model} & \textbf{Wavelength} & \textbf{Primary Observables} \\
\midrule
Ceilometer & Vaisala CL-31 & 905\,nm & $\beta_\mathrm{att}(z,t)$, cloud base height \\
Doppler Lidar & Halo Photonics StreamLine XR+ & 1550\,nm & $w(z,t)$, $\mathbf{U}(z,t)$, $\beta_\mathrm{att}$ \\
MWR & Radiometrics MP-3000A & 22--58\,GHz & $T(z,t)$, $q_v(z,t)$, LWP, PWV \\
\bottomrule
\end{tabular}
\end{center}

\subsection{Microphysics Compatibility}

The backscatter forward model adapts automatically to the active microphysics scheme:

\begin{itemize}
  \item \textbf{Bulk schemes (Kessler, Morrison, SAM, SatAdj, WSM6, Null)} ---
        Use a fixed effective radius $r_\mathrm{eff}$ with bulk LWC/IWC extracted from
        \texttt{RhoQ2\_comp} (liquid) and \texttt{RhoQ3\_comp} (ice) in the conserved-state array
        (Section~\ref{sec:ceil_simple}).
  \item \textbf{SBM} (Spectral Bin Microphysics) --- Eulerian bin-resolving scheme; backscatter
        uses per-bin number concentration $N_b$ [$\mathrm{m}^{-3}$] and bin-centre radius $r_b$,
        identical to the \texttt{sbm\_N\_cloud} and \texttt{sbm\_reflectivity} plotfile
        variables (Section~\ref{sec:size_resolved}).
  \item \textbf{SDM} (Super-Droplet Method) --- Lagrangian particle-based scheme; backscatter
        uses each super-droplet's physical radius $r_p$ and multiplicity $\xi_p$
        (Section~\ref{sec:size_resolved}).
\end{itemize}

The MWR simulator uses only bulk thermodynamic variables and behaves identically
for all microphysics schemes.

% ============================================================
\section{Ceilometer Simulator (Vaisala CL-31, 905\,nm)}
% ============================================================

\subsection{Instrument Description}

The Vaisala CL-31 is a single-lens backscatter lidar operating at $\lambda = 905$\,nm.
It transmits short laser pulses and measures time-gated return signals.
The primary observable is the \textbf{attenuated backscatter} coefficient
$\beta_\mathrm{att}(z,t)$ [$\mathrm{m}^{-1}\,\mathrm{sr}^{-1}$] as a function of height
and time.

\subsection{Simple Backscatter Forward Model (bulk schemes)}\label{sec:ceil_simple}

\subsubsection{Volume Extinction Coefficient}

For cloud droplets with effective radius $r_\mathrm{eff} \gg \lambda$ (geometric optics regime,
valid for $r_\mathrm{eff} \gtrsim 5\,\mu\mathrm{m}$):

\begin{equation}
  \sigma_\mathrm{ext}(z) = \frac{3}{2}\,\frac{\mathrm{LWC}(z)}{\rho_w\,r_\mathrm{eff}(z)}
  \quad [\mathrm{m}^{-1}]
  \label{eq:sigma_ext_simple}
\end{equation}

\noindent where:
\begin{align*}
  \mathrm{LWC}(z) &= \rho_\mathrm{air}(z)\,q_c(z) \quad [\mathrm{kg\,m}^{-3}] \\
  \rho_w &= 1000\;\mathrm{kg\,m}^{-3} \quad \text{(liquid water density)} \\
  r_\mathrm{eff}(z) &= \frac{\int_0^\infty n(r)\,r^3\,dr}{\int_0^\infty n(r)\,r^2\,dr}
  \quad \text{(third/second moment of size distribution)}
\end{align*}

For ice clouds the same form applies using IWC and $r_{\mathrm{eff},i}$.
In the absence of size-distribution information (all bulk schemes),
the defaults are $r_\mathrm{eff} = 10\,\mu\mathrm{m}$ (liquid) and
$r_{\mathrm{eff},i} = 30\,\mu\mathrm{m}$ (ice), both user-configurable.

Both liquid and ice contributions to extinction and backscatter are summed before
applying the two-way optical depth:
\begin{equation}
  \sigma_\mathrm{tot}(z) = \sigma_\mathrm{liq}(z) + \sigma_\mathrm{ice}(z)
\end{equation}

\subsubsection{Volume Backscatter Coefficient}

The lidar ratio $S_L$ relates backscatter to extinction:
\begin{equation}
  \beta(z) = \frac{\sigma_\mathrm{ext}(z)}{S_L}
  \quad [\mathrm{m}^{-1}\,\mathrm{sr}^{-1}]
  \label{eq:beta_simple}
\end{equation}

\noindent Adopted lidar ratio values:
\begin{center}
\begin{tabular}{lll}
\toprule
Target & Wavelength & $S_L$ \\
\midrule
Liquid cloud & 905\,nm  & 18.8\,sr \\
Ice cloud    & 905\,nm  & 25.0\,sr \\
Liquid cloud & 1550\,nm & 20.0\,sr \\
Ice cloud    & 1550\,nm & 28.0\,sr \\
\bottomrule
\end{tabular}
\end{center}

\subsubsection{Two-Way Optical Depth and Transmission}

The cumulative two-way optical depth from the surface to height $z$:
\begin{equation}
  \tau(z) = 2\int_0^z \sigma_\mathrm{ext}(z')\,dz'
  \label{eq:tau}
\end{equation}

Discretized (trapezoidal rule, $k$ = vertical grid index):
\begin{equation}
  \tau_k = \tau_{k-1} + \bigl(\sigma_{\mathrm{tot},k} + \sigma_{\mathrm{tot},k-1}\bigr)\,\Delta z_k
\end{equation}

\subsubsection{Attenuated Backscatter}

\begin{equation}
  \boxed{\beta_\mathrm{att}(z) = \beta(z)\,\exp\bigl[-\tau(z)\bigr]}
  \quad [\mathrm{m}^{-1}\,\mathrm{sr}^{-1}]
  \label{eq:beta_att}
\end{equation}

This is the quantity directly observed by the ceilometer.

\subsection{Size-Resolved Geometric-Optics Backscatter (SBM and SDM)}\label{sec:size_resolved}

When SBM or SDM microphysics is active, the extinction coefficient is computed from the
actual droplet/crystal size distribution using the geometric-optics limit ($Q_\mathrm{ext} = 2$,
valid for $r \gg \lambda$, i.e.\ cloud droplets $r \gtrsim 5\,\mu\mathrm{m}$).
The same variables and bin radii are used as in the ERF plotfile diagnostics
(\texttt{sbm\_N\_cloud}, \texttt{sbm\_reflectivity}), ensuring consistency between
plotfile and instrument-simulator outputs.

\subsubsection{SBM Path}

For each liquid (or ice) bin $b$ with bin-centre radius $r_b$ [$\mathrm{m}$] and number
concentration $N_b$ [$\mathrm{m}^{-3}$]:
\begin{equation}
  \sigma_{\mathrm{ext},b}(z) = N_b(z)\cdot 2\pi r_b^2
  \quad [\mathrm{m}^{-1}]
  \label{eq:sigma_ext_sbm}
\end{equation}
\noindent where $N_b$ is read from
\texttt{cons[RhoQ1\_comp + idx\_nc\_start + b]} (two-moment SBM, $N_b$ in $\mathrm{m}^{-3}$),
and $r_b = \texttt{GetDiameterGridHostUM()}[b] / 2$.
The column extinction is the sum over all 33 liquid (or ice) bins:
\begin{equation}
  \sigma_\mathrm{ext,liq}(z) = \sum_{b=0}^{32}\sigma_{\mathrm{ext},b}^{\mathrm{liq}}(z),
  \qquad
  \sigma_\mathrm{ext,ice}(z) = \sum_{b=0}^{32}\sigma_{\mathrm{ext},b}^{\mathrm{ice}}(z)
\end{equation}
A one-moment fallback derives $N_b$ from the mass density bin if
\texttt{idx\_nc\_start} $< 0$.

\subsubsection{SDM Path}

For each super-droplet $p$ in a column $(i_\mathrm{loc}, j_\mathrm{loc})$, its contribution
to vertical cell $k$ is accumulated on the GPU using atomic addition:
\begin{equation}
  \sigma_\mathrm{ext}(k) \mathrel{+}=
  \frac{\xi_p \cdot 2\pi r_p^2}{V_\mathrm{cell}}
  \quad [\mathrm{m}^{-1}]
  \label{eq:sigma_ext_sdm}
\end{equation}
\noindent where $\xi_p$ is the multiplicity (real droplet count represented),
$r_p$ [$\mathrm{m}$] is the physical droplet radius, and
$V_\mathrm{cell} = \Delta x\,\Delta y\,\Delta z$ is the grid cell volume.
The vertical index $k$ is obtained from \texttt{p.idata(ERFParticlesIntIdxAoS::k)}.
Accumulation uses \texttt{amrex::Gpu::Atomic::AddNoRet} on a \texttt{Gpu::DeviceVector},
followed by device-to-host copy and \texttt{MPI\_Reduce} to the IO processor.
SDM does not separately tag liquid vs.\ ice phase; the full $\sigma_\mathrm{ext}$
is stored in \texttt{sigma\_liq}, with \texttt{sigma\_ice} = 0.

\subsubsection{Data Pathway}

The pre-computed $\sigma$ columns are packaged into an \texttt{InstrMicrophySizeData}
struct and passed to \texttt{write()} as an optional pointer:
\begin{lstlisting}[language=c++]
struct InstrMicrophySizeData {
    Vector<Vector<Real>> sigma_liq;  // [site][k]  m-1
    Vector<Vector<Real>> sigma_ice;  // [site][k]  m-1
    bool empty() const { return sigma_liq.empty(); }
};
// In ERF::Evolve() -- only on write steps (avoids overhead every step):
ceil_mpdata = fill_sbm(ilocs, jlocs);   // or fill_sdm(...)
m_ceilometer_sim->write(..., ceil_mpdata.get());
\end{lstlisting}
Backscatter is computed via
\texttt{InstrumentSimUtil::compute\_beta\_att\_from\_sigma()},
which applies the optical-depth integral and lidar-ratio conversion
(equations~\ref{eq:tau}--\ref{eq:beta_att}).
When \texttt{mpdata == nullptr} (any bulk scheme), the simulator falls back to
the simple fixed-$r_\mathrm{eff}$ model (Section~\ref{sec:ceil_simple}).

% ============================================================
\subsection{Mie Scattering Forward Model (Planned --- Not Yet Implemented)}\label{sec:mie}

\textit{Note: The full Mie backscatter series is not yet implemented.
When \texttt{backscatter\_model = mie} is specified, the code currently falls back to
the size-resolved geometric-optics model of Section~\ref{sec:size_resolved}
(SBM/SDM) or the bulk fixed-$r_\mathrm{eff}$ model (bulk schemes).
The section below documents the planned Mie algorithm for future implementation.}

\subsubsection{Mie Efficiency Factors}

For each SBM bin $n$ with bin-center radius $r_n$ at wavelength $\lambda$, the size parameter is:
\begin{equation}
  x_n = \frac{2\pi r_n}{\lambda}
\end{equation}

The Mie extinction and backscatter efficiency factors $Q_\mathrm{ext}(x_n,m)$ and
$Q_\mathrm{back}(x_n,m)$ are computed via the Bohren \& Huffman (1983) Mie series
\cite{BohrenHuffman1983}:

\begin{align}
  Q_\mathrm{ext} &= \frac{2}{x^2}\sum_{l=1}^{N}(2l+1)\,\mathrm{Re}(a_l + b_l)
  \label{eq:Qext_Mie} \\[6pt]
  Q_\mathrm{back} &= \frac{1}{x^2}\left|\sum_{l=1}^{N}(2l+1)(-1)^l(a_l - b_l)\right|^2
  \label{eq:Qback_Mie}
\end{align}

\noindent where $a_l$, $b_l$ are the Mie coefficients (complex Riccati-Bessel functions),
$m = n_r - i n_i$ is the complex refractive index of water/ice.

Refractive indices used:
\begin{itemize}
  \item Liquid water at 905\,nm: $m = 1.328 - 1.67\times10^{-8}i$
  \item Liquid water at 1550\,nm: $m = 1.318 - 1.18\times10^{-4}i$
  \item Ice at 905\,nm: $m = 1.304 - 3.65\times10^{-9}i$
\end{itemize}

These tables are precomputed at initialization using the SBM bin diameters
from \texttt{ERF\_BinGrid.H} (\texttt{m\_diameter\_grid\_host}).

\subsubsection{Column-Integrated Mie Backscatter}

\begin{align}
  \beta_\mathrm{Mie}(z) &= \sum_n N_n(z)\,\pi r_n^2\,Q_{\mathrm{back},n}
  \quad [\mathrm{m}^{-1}\,\mathrm{sr}^{-1}]
  \label{eq:beta_Mie} \\[4pt]
  \sigma_\mathrm{Mie}(z) &= \sum_n N_n(z)\,\pi r_n^2\,Q_{\mathrm{ext},n}
  \quad [\mathrm{m}^{-1}]
  \label{eq:sigma_Mie}
\end{align}

\noindent where $N_n(z)$ [$\mathrm{m}^{-3}$] is the number concentration in bin $n$ at height $z$.
The attenuated backscatter is then:
\begin{equation}
  \beta_\mathrm{att}(z) = \beta_\mathrm{Mie}(z)\,\exp\!\left(-2\int_0^z \sigma_\mathrm{Mie}(z')\,dz'\right)
\end{equation}

\subsubsection{SDM (Lagrangian) Adaptation}

For SDM, at each grid cell $(i,j,k)$ the super-droplet radii are histogrammed into
a temporary array matching the SBM bin grid. The number concentration per bin is estimated as:
\begin{equation}
  N_n = \frac{1}{V_\mathrm{cell}}\sum_{p\,:\,r_p \in [r_n^-, r_n^+]} \xi_p
\end{equation}
where $\xi_p$ is the multiplicity (real droplet count) of super-droplet $p$,
and $V_\mathrm{cell}$ is the grid cell volume.

\subsection{Cloud Base Detection}

Cloud base height is the lowest level where $\beta_\mathrm{att}$ exceeds a threshold:
\begin{equation}
  z_\mathrm{cb} = \min\{z : \beta_\mathrm{att}(z) > \beta_\mathrm{thresh}\}
\end{equation}

Default: $\beta_\mathrm{thresh} = 3\times10^{-4}\,\mathrm{m}^{-1}\,\mathrm{sr}^{-1}$.
Output is $-9999$ if no cloud is detected.

% ============================================================
\section{Doppler Lidar Simulator (Halo StreamLine XR+, 1550\,nm)}
% ============================================================

\subsection{Instrument Description}

The Halo Photonics StreamLine XR+ is a coherent pulsed Doppler lidar at $\lambda=1550$\,nm.
It operates in multiple scan modes:
\begin{itemize}
  \item \textbf{Vertical stare}: beam pointing zenith, measures $w(z,t)$ and CNR$(z,t)$
  \item \textbf{VAD scan}: conical scan at fixed elevation angle, retrieves $u(z,t)$, $v(z,t)$
  \item \textbf{Attenuated backscatter}: passive measurement of $\beta_\mathrm{att}(z,t)$
\end{itemize}
Each mode is controlled independently via input parameters.

\subsection{Vertical Stare Mode}

For a zenith-pointing beam, the radial (Doppler) velocity equals the vertical wind component:
\begin{equation}
  V_r(z) = w(z) \quad [\mathrm{m\,s}^{-1}]
\end{equation}

The cell-centred vertical velocity is the average of the two enclosing face values:
$w_k = \tfrac{1}{2}(w_{k-1/2} + w_{k+1/2})$.

\subsubsection{CNR Model Selection (\texttt{cnr\_model})}

The Carrier-to-Noise Ratio (CNR) is output alongside $w(z,t)$ in the vertical-stare file.
Two models are available.

\paragraph{Normalized model (default, \texttt{cnr\_model = normalized})}

A dimensionless range-normalized proxy convenient for qualitative comparison:
\begin{equation}
  \mathrm{CNR}(r) = \frac{\beta_\mathrm{att}(r)}{r^2}
  \label{eq:cnr_normalized}
\end{equation}
This requires \texttt{do\_backscatter = true} to supply $\beta_\mathrm{att}$;
otherwise CNR is output as zero.

\paragraph{Physical model (\texttt{cnr\_model = physical})}

Based on the coherent Doppler lidar performance equation of Targ et al.\ (1991)
\cite{Targ1991}.  The CNR is related to the attenuated backscatter through the
instrument function $F(r)$:

\begin{equation}
  \boxed{\mathrm{CNR}(r) = F(r)\;\beta_\mathrm{att}(r)}
  \label{eq:cnr_physical}
\end{equation}

where
\begin{equation}
  F(r) = \frac{E\,\eta\,\lambda\,\pi\,\omega^2}{2\,h\,B_w\,r^2}
  \cdot\underbrace{\left[1 + \left(1-\frac{r}{R}\right)^2
    \!\left(\frac{\pi\omega^2}{\lambda r}\right)^{\!2}\right]^{-1}}_{\text{Gaussian beam focus correction}}
  \label{eq:F_targ}
\end{equation}

\noindent The parameters are:
\begin{center}
\begin{tabular}{llll}
\toprule
Symbol    & Parameter name            & Meaning                                   & Default (Halo XR+) \\
\midrule
$E$       & \texttt{pulse\_energy}    & Pulse energy [J]                          & $1.5\times10^{-6}$\,J \\
$\eta$    & \texttt{optics\_efficiency}& Optical + detector efficiency [-]        & 0.1 \\
$\lambda$ & ---                       & Wavelength [m] (fixed)                    & $1.55\times10^{-6}$\,m \\
$\omega$  & \texttt{beam\_radius}     & $1/e^2$ beam radius at aperture [m]       & 0.0225\,m \\
$h$       & ---                       & Planck constant [J\,s] (fixed)            & $6.626\times10^{-34}$\,J\,s \\
$B_w$     & \texttt{bandwidth}        & Receiver noise bandwidth [Hz]             & $50\times10^6$\,Hz \\
$R$       & \texttt{focus\_range}     & Radius of curvature at aperture [m]       & $10^9$\,m (collimated) \\
\bottomrule
\end{tabular}
\end{center}

The factor $\pi\omega^2/(\lambda r)$ is the Rayleigh-range-to-range ratio.
The focus correction equals 1 when $r = R$ (beam focused at range $R$) and
approaches 1 everywhere for a collimated beam ($R \to \infty$).

\subsection{VAD (Velocity-Azimuth Display) Mode}

\subsubsection{Radial Velocity Simulation}

For a beam at elevation angle $\phi$ from horizontal and azimuth $\psi$ (clockwise from North),
the beam unit vector in East-North-Up coordinates is:
\begin{equation}
  \hat{d} = (\cos\phi\cos\psi,\;\cos\phi\sin\psi,\;\sin\phi)
\end{equation}

The simulated radial velocity (projection of wind onto beam):
\begin{equation}
  \boxed{V_r(\psi, z) = u(z)\cos\phi\cos\psi + v(z)\cos\phi\sin\psi + w(z)\sin\phi}
  \label{eq:vad_radial}
\end{equation}

Sampled at $N_\psi = 36$ azimuths: $\psi_m = m \times 10^\circ$, $m = 0, 1, \ldots, 35$.
Cell-centred $u$ and $v$ are computed by averaging the two enclosing face values.

\subsubsection{VAD Retrieval}

The radial velocity field is modeled as a truncated Fourier series in azimuth:
\begin{equation}
  V_r(\psi, z) \approx A_0(z) + A_1(z)\cos\psi + B_1(z)\sin\psi
  \label{eq:vad_fit}
\end{equation}

Coefficients are determined by least-squares fit over the $N_\psi$ azimuths.
The horizontal wind components are retrieved as:
\begin{align}
  w(z) &= A_0(z)/\sin\phi \label{eq:vad_w} \\
  u(z) &= A_1(z)/\cos\phi \label{eq:vad_u} \\
  v(z) &= B_1(z)/\cos\phi \label{eq:vad_v}
\end{align}

Horizontal wind speed and meteorological wind direction:
\begin{align}
  |V_h|(z) &= \sqrt{u(z)^2 + v(z)^2} \\
  \mathrm{WD}(z) &= \mathrm{atan2}(-u,-v)\times\frac{180}{\pi} + 180
  \quad [\text{degrees, meteorological convention}]
\end{align}

\subsection{Attenuated Backscatter at 1550\,nm}

Same forward model as the ceilometer (\autoref{eq:sigma_ext_simple}--\ref{eq:beta_att})
but with $\lambda = 1550$\,nm lidar ratios ($S_L = 20$\,sr for liquid, $28$\,sr for ice).
Both liquid and ice cloud contributions are included:
\begin{equation}
  \beta_\mathrm{att}(z) = \bigl[\beta_\mathrm{liq}(z) + \beta_\mathrm{ice}(z)\bigr]\,
  \exp\!\left(-\tau_\mathrm{liq}(z) - \tau_\mathrm{ice}(z)\right)
\end{equation}

% ============================================================
\section{Microwave Radiometer Simulator (Radiometrics MP-3000A)}
% ============================================================

\subsection{Instrument Description}

The Radiometrics MP-3000A is a passive microwave radiometer with 35 channels:
\begin{itemize}
  \item \textbf{K-band}: 22 channels from 22.235 to 40.0\,GHz --- water vapor absorption
  \item \textbf{V-band}: 13 channels from 51.25 to 58.0\,GHz --- oxygen absorption complex
\end{itemize}

The instrument retrieves temperature profiles, water vapor profiles, LWP, and PWV using
neural network algorithms applied to observed brightness temperatures.

\subsection{Direct Profile Output}

\subsubsection{Temperature Profile}

Model potential temperature $\theta$ is converted to absolute temperature via the Exner function:
\begin{align}
  T(z) &= \theta(z)\,\Pi(z) \label{eq:exner_T} \\
  \Pi(z) &= \left(\frac{p(z)}{p_0}\right)^{R_d/c_p} \label{eq:exner}
\end{align}

\noindent where:
\begin{align*}
  \theta(z) &= (\rho\theta)(z) / \rho(z) \quad\text{from state array}\\
  p(z) &= p_0(z) \quad\text{(base-state pressure, 1D column)}\\
  p_0 &= 10^5\;\mathrm{Pa} \quad\text{(reference pressure)}\\
  R_d &= 287.04\;\mathrm{J\,kg}^{-1}\mathrm{K}^{-1} \\
  c_p &= 1004.5\;\mathrm{J\,kg}^{-1}\mathrm{K}^{-1}
\end{align*}

\subsubsection{Water Vapor and Cloud Liquid Profiles}

\begin{align}
  q_v(z) &= \frac{(\rho q_v)(z)}{\rho(z)}
  = \frac{\mathtt{state}(i,j,k,\,\text{RhoQ1\_comp})}{\mathtt{state}(i,j,k,\,\text{Rho\_comp})}
  \quad [\mathrm{kg\,kg}^{-1}] \\[4pt]
  q_c(z) &= \frac{(\rho q_c)(z)}{\rho(z)}
  = \frac{\mathtt{state}(i,j,k,\,\text{RhoQ2\_comp})}{\mathtt{state}(i,j,k,\,\text{Rho\_comp})}
  \quad [\mathrm{kg\,kg}^{-1}]
\end{align}

Both $q_v$ and $q_c$ are output in $\mathrm{g\,kg}^{-1}$ in the profile file.

\subsubsection{Column Integrals}

Liquid Water Path (LWP):
\begin{equation}
  \mathrm{LWP} = \int_0^{z_\mathrm{top}} \rho(z)\,q_c(z)\,dz
  \approx \sum_k \rho_k\,q_{c,k}\,\Delta z_k
  \quad [\mathrm{kg\,m}^{-2}]
  \label{eq:lwp}
\end{equation}
Output in $\mathrm{g\,m}^{-2}$ (multiply by 1000).

Precipitable Water Vapor (PWV):
\begin{equation}
  \mathrm{PWV} = \frac{1}{\rho_w}\int_0^{z_\mathrm{top}} \rho(z)\,q_v(z)\,dz
  \approx \frac{1}{\rho_w}\sum_k \rho_k\,q_{v,k}\,\Delta z_k
  \quad [\mathrm{mm}]
  \label{eq:pwv}
\end{equation}

\noindent where $\rho_w = 1000\;\mathrm{kg\,m}^{-3}$ and the result in mm equals
the column integral in $\mathrm{kg\,m}^{-2}$ (numerically).

\subsection{Brightness Temperature Forward Model (Optional)}

\subsubsection{Gaseous Absorption}

The total volume absorption coefficient at frequency $f$ [GHz] and height $z$ is:
\begin{equation}
  \alpha(f, z) = \alpha_{\mathrm{H_2O}}(f, T, p, \rho_v)
               + \alpha_{\mathrm{O_2}}(f, T, p)
               + \alpha_{\mathrm{liq}}(f, T, \mathrm{LWC})
  \quad [\mathrm{Np\,m}^{-1}]
\end{equation}

\paragraph{Water vapor absorption (K-band)}
Simplified Rosenkranz (1998) model \cite{Rosenkranz1998} with dominant line at 22.235\,GHz
plus a continuum term.  The Van Vleck-Weisskopf lineshape is:
\begin{equation}
  F_i(f) = \frac{f}{f_0}\left[\frac{\gamma}{(f-f_0)^2+\gamma^2}
                               + \frac{\gamma}{(f+f_0)^2+\gamma^2}\right]
\end{equation}
Line width $\gamma$ depends on pressure and water vapor density following
Rosenkranz (1998) Table~1.

\paragraph{Oxygen absorption (V-band)}
Simplified Liebe et al.\ (1993) \cite{Liebe1993} oxygen complex.
Thirty dominant lines from 50.474 to 66.302\,GHz are summed with pressure-broadened
Van Vleck-Weisskopf profiles and line strengths scaled as $\theta^3\,p$.

\paragraph{Liquid water absorption}
Rayleigh approximation valid for $f \ll 300$\,GHz:
\begin{equation}
  \alpha_\mathrm{liq}(f) = \frac{6\pi f}{c}\,
  \frac{3\,\varepsilon''}{(\varepsilon_r + 2)^2 + (\varepsilon'')^2}\,
  \frac{\mathrm{LWC}}{\rho_w}
  \label{eq:alpha_liq}
\end{equation}

\noindent The complex permittivity $\varepsilon = \varepsilon_r - i\varepsilon''$ of liquid water
uses the Ellison (2007) double-Debye parameterization \cite{Ellison2007}.

\subsubsection{Radiative Transfer}

For zenith-pointing geometry, the upwelling brightness temperature is:
\begin{equation}
  \boxed{T_B(f) = T_\mathrm{bg}\,e^{-\tau_\mathrm{tot}(f)}
  + \int_0^\infty T(z)\,\alpha(f,z)\,
  e^{-\tau(f,\,0,\,z)}\,dz}
  \label{eq:RT}
\end{equation}

\noindent where:
\begin{align*}
  T_\mathrm{bg} &= 2.75\;\mathrm{K} \quad\text{(cosmic microwave background)} \\
  \tau(f, z_1, z_2) &= \int_{z_1}^{z_2}\alpha(f,z')\,dz' \\
  \tau_\mathrm{tot}(f) &= \tau(f, 0, \infty)
\end{align*}

Discrete form (summed over $K$ vertical levels, $\tau_k$ = optical depth from surface to $z_k$):
\begin{equation}
  T_B(f) \approx T_\mathrm{bg}\,e^{-\tau_K}
  + \sum_{k=1}^{K} T_k\,\bigl(e^{-\tau_{k-1}} - e^{-\tau_k}\bigr)
\end{equation}

\subsubsection{MP-3000A Channel Frequencies}

\begin{center}
\begin{tabular}{rl|rl}
\toprule
\multicolumn{2}{c|}{\textbf{K-band (22 channels)}} &
\multicolumn{2}{c}{\textbf{V-band (13 channels)}} \\
Ch & $f$ [GHz] & Ch & $f$ [GHz] \\
\midrule
1  & 22.235 & 23 & 51.250 \\
2  & 22.500 & 24 & 51.760 \\
3  & 23.035 & 25 & 52.280 \\
4  & 23.835 & 26 & 52.800 \\
5  & 25.000 & 27 & 53.340 \\
6  & 26.235 & 28 & 53.860 \\
7  & 28.000 & 29 & 54.400 \\
8  & 30.000 & 30 & 54.940 \\
9  & 31.400 & 31 & 55.500 \\
10 & 32.000 & 32 & 56.020 \\
11 & 32.400 & 33 & 56.660 \\
12 & 33.000 & 34 & 57.290 \\
13 & 34.000 & 35 & 58.000 \\
14 & 35.000 & & \\
15 & 36.000 & & \\
16 & 36.500 & & \\
17 & 37.000 & & \\
18 & 37.500 & & \\
19 & 38.000 & & \\
20 & 38.600 & & \\
21 & 39.500 & & \\
22 & 40.000 & & \\
\bottomrule
\end{tabular}
\end{center}

% ============================================================
\section{Input File Parameters}
% ============================================================

All parameters use the \texttt{erf.*} ParmParse prefix.
Multiple instrument sites can be specified by providing equal-length lists for
\texttt{i\_loc} and \texttt{j\_loc}; indices are 0-based cell indices.

\begin{lstlisting}[language=bash,caption={Ceilometer input parameters (defaults shown)}]
# Ceilometer (Vaisala CL-31, 905 nm)
erf.ceilometer.do_sim             = true
erf.ceilometer.i_loc              = 32 48    # REQUIRED; 0-based i indices (list OK)
erf.ceilometer.j_loc              = 32 48    # REQUIRED; must match i_loc count
erf.ceilometer.output_interval    = 10       # write every N timesteps
erf.ceilometer.output_file        = ceil_    # prefix; ceil_00.txt, ceil_01.txt, ...
erf.ceilometer.backscatter_model  = simple   # "simple" (geometric optics) or "mie"
erf.ceilometer.r_eff_default      = 1.0e-5  # m; liquid effective radius fallback
erf.ceilometer.r_eff_ice_default  = 3.0e-5  # m; ice effective radius fallback
erf.ceilometer.beta_thresh        = 3.0e-4  # m-1 sr-1; cloud base detection threshold
\end{lstlisting}

\begin{lstlisting}[language=bash,caption={Doppler Lidar input parameters (defaults shown)}]
# Doppler Lidar (Halo StreamLine XR+, 1550 nm)
erf.doppler_lidar.do_sim              = true
erf.doppler_lidar.i_loc               = 32      # REQUIRED; one or more sites
erf.doppler_lidar.j_loc               = 32      # REQUIRED
erf.doppler_lidar.output_interval     = 10
erf.doppler_lidar.output_file         = dl_
erf.doppler_lidar.backscatter_model   = simple  # "simple" or "mie"
erf.doppler_lidar.r_eff_default       = 1.0e-5  # m; liquid effective radius fallback
erf.doppler_lidar.do_vertical_stare   = true    # output w(z,t) and CNR  [default: true]
erf.doppler_lidar.do_vad              = false   # output u,v,wind speed/dir [default: false]
erf.doppler_lidar.vad_elevation_angle = 75.0    # degrees from horizontal
erf.doppler_lidar.do_backscatter      = false   # output beta_att(z,t)  [default: false]
# CNR model for vertical stare output:
erf.doppler_lidar.cnr_model           = normalized  # "normalized" (default) or "physical"
# Physical CNR parameters (Targ 1991) -- only used when cnr_model = physical:
erf.doppler_lidar.pulse_energy        = 1.5e-6  # E   [J]
erf.doppler_lidar.optics_efficiency   = 0.1     # eta [-]
erf.doppler_lidar.beam_radius         = 0.0225  # omega [m]  1/e^2 beam radius
erf.doppler_lidar.bandwidth           = 50.0e6  # Bw  [Hz]  receiver bandwidth
erf.doppler_lidar.focus_range         = 1.0e9   # R   [m]   (1e9 = collimated)
\end{lstlisting}

\begin{lstlisting}[language=bash,caption={MWR input parameters (defaults shown)}]
# MWR (Radiometrics MP-3000A, 35 channels)
erf.mwr.do_sim             = true
erf.mwr.i_loc              = 32      # REQUIRED
erf.mwr.j_loc              = 32      # REQUIRED
erf.mwr.output_interval    = 10
erf.mwr.output_file        = mwr_
erf.mwr.do_brightness_temp = false   # true = also write simulated TB (35 channels)
\end{lstlisting}

% ============================================================
\section{Output File Formats}
% ============================================================

All output files are ASCII text with comment headers (\texttt{\#} prefix lines).
Time rows are appended at each write interval.

\subsection{Ceilometer: \texttt{ceil\_XX.txt}}
\begin{lstlisting}
# ERF Ceilometer Simulator (Vaisala CL-31, 905 nm)
# i_loc=32  j_loc=32
# backscatter_model=simple
# Col 1: time [s]
# Col 2: height [m]
# Col 3: att_backscatter [m-1 sr-1]
# Col 4: cloud_base_height [m]  (-9999 = no cloud)
  1.000000e+01  2.500000e+01  1.234000e-06  2.500000e+02
  1.000000e+01  7.500000e+01  3.456000e-04  2.500000e+02
\end{lstlisting}

\subsection{Doppler Lidar Stare: \texttt{dl\_XX\_stare.txt}}
\begin{lstlisting}
# ERF Doppler Lidar Simulator (Halo XR+, 1550 nm) -- vertical stare
# i_loc=32  j_loc=32
# cnr_model=normalized  (CNR = beta_att / r^2)
# Col 1: time [s],  Col 2: height [m]
# Col 3: w [m s-1],  Col 4: CNR [-]
  1.000000e+01  2.500000e+01  1.230000e-01  4.560000e-03
\end{lstlisting}

\subsection{Doppler Lidar VAD: \texttt{dl\_XX\_vad.txt}}
\begin{lstlisting}
# ERF Doppler Lidar Simulator -- VAD mode (elev=75.0 deg, 36 azimuths)
# Col 1: time [s],  Col 2: height [m]
# Col 3: u [m s-1],  Col 4: v [m s-1]
# Col 5: wind_speed [m s-1],  Col 6: wind_dir [deg, met convention]
\end{lstlisting}

\subsection{Doppler Lidar Backscatter: \texttt{dl\_XX\_beta.txt}}
\begin{lstlisting}
# ERF Doppler Lidar Simulator -- att. backscatter (1550 nm)
# Col 1: time [s],  Col 2: height [m],  Col 3: beta_att [m-1 sr-1]
\end{lstlisting}

\subsection{MWR Profiles: \texttt{mwr\_XX\_profiles.txt}}
\begin{lstlisting}
# ERF MWR Simulator (Radiometrics MP-3000A)
# Col 1: time [s],  Col 2: height [m]
# Col 3: T [K],  Col 4: qv [g kg-1],  Col 5: qc [g kg-1]
\end{lstlisting}

\subsection{MWR Column Integrals: \texttt{mwr\_XX\_column.txt}}
\begin{lstlisting}
# ERF MWR Simulator -- column integrals
# Col 1: time [s],  Col 2: LWP [g m-2],  Col 3: PWV [mm]
  1.000000e+01  2.534000e+01  1.872000e+01
\end{lstlisting}

\subsection{MWR Brightness Temperatures: \texttt{mwr\_XX\_tb.txt} (optional)}
\begin{lstlisting}
# ERF MWR Simulator -- simulated brightness temperatures (35 channels)
# Col 1: time [s],  Col 2: frequency [GHz],  Col 3: TB [K]
  1.000000e+01  2.223500e+01  2.654300e+02
  1.000000e+01  2.250000e+01  2.671200e+02
  ...  (35 rows per timestep)
\end{lstlisting}

% ============================================================
\section{Software Architecture}
% ============================================================

\subsection{Source File}

\begin{center}
\begin{tabular}{ll}
\toprule
File & Purpose \\
\midrule
\texttt{Source/IO/ERF\_InstrumentSimulators.H} & Header-only: all three simulator structs \\
\bottomrule
\end{tabular}
\end{center}

\subsection{Modified Files}

\begin{center}
\begin{tabular}{ll}
\toprule
File & Change \\
\midrule
\texttt{Source/ERF.H} & Add \texttt{unique\_ptr} members + method declarations \\
\texttt{Source/ERF.cpp} & Initialize in \texttt{InitData()}, call in \texttt{Evolve()} and \texttt{post\_timestep()} \\
\texttt{Source/IO/Make.package} & Add \texttt{ERF\_InstrumentSimulators.H} to \texttt{CEXE\_headers} \\
\texttt{CMake/BuildERFExe.cmake} & Add \texttt{ERF\_InstrumentSimulators.H} to target sources \\
\bottomrule
\end{tabular}
\end{center}

\subsection{Design Pattern}

The simulator structs follow the \texttt{LineSampler} / \texttt{PlaneSampler} pattern in
\texttt{Source/IO/ERF\_SampleData.H}:
\begin{itemize}
  \item \texttt{init()} reads \texttt{ParmParse} parameters and opens output files.
  \item \texttt{write()} checks interval, extracts z-columns from MultiFabs at the
        requested $(i,j)$ location(s), computes derived quantities, and writes ASCII output.
  \item Column extraction is GPU-safe: \texttt{InstrumentSimUtil::extract\_column()} and
        \texttt{extract\_face\_column()} use \texttt{amrex::ParallelFor} on the device
        followed by device-to-host copy and \texttt{MPI\_Reduce} to the IO processor.
  \item All file I/O is guarded by \texttt{ParallelDescriptor::IOProcessor()}.
  \item The optional \texttt{InstrMicrophySizeData*} argument to the ceilometer and
        Doppler lidar \texttt{write()} functions carries pre-computed SBM or SDM
        extinction columns; \texttt{nullptr} invokes the bulk fixed-$r_\mathrm{eff}$ path.
\end{itemize}

\subsection{Variable Access Summary}

\begin{center}
\begin{tabular}{lll}
\toprule
Quantity & ERF Access & Notes \\
\midrule
$\rho$ & \texttt{cons(i,j,k,Rho\_comp)} & \\
$\rho\theta$ & \texttt{cons(i,j,k,RhoTheta\_comp)} & \\
$\rho q_v$ & \texttt{cons(i,j,k,RhoQ1\_comp)} & \\
$\rho q_c$ & \texttt{cons(i,j,k,RhoQ2\_comp)} & bulk liquid (all schemes) \\
$\rho q_i$ & \texttt{cons(i,j,k,RhoQ3\_comp)} & bulk ice (guarded by \texttt{nComp()}) \\
$u$ & \texttt{xvel(i,j,k)} & face-centred; averaged for cell-centre \\
$v$ & \texttt{yvel(i,j,k)} & face-centred; averaged for cell-centre \\
$w$ & \texttt{zvel(i,j,k)} & face-centred; averaged for cell-centre \\
$p_0(z)$ & \texttt{base\_state(i,j,k,BaseState::p0\_comp)} & 1D base pressure \\
SBM $N_b$ & \texttt{cons[RhoQ1\_comp + idx\_nc\_start + b]} & two-moment, $[\mathrm{m}^{-3}]$ \\
SBM $r_b$ & \texttt{GetDiameterGridHostUM()[b] / 2} & precomputed table, $[\mathrm{m}]$ \\
SDM $r_p$ & \texttt{soa.GetRealData(rtoff\_r + radius)[p]} & per-particle radius $[\mathrm{m}]$ \\
SDM $\xi_p$ & \texttt{soa.GetRealData(rtoff\_r + multiplicity)[p]} & multiplicity $[-]$ \\
SDM $k_p$ & \texttt{p.idata(ERFParticlesIntIdxAoS::k)} & vertical cell index \\
\bottomrule
\end{tabular}
\end{center}

% ============================================================
\section{Verification and Testing}
% ============================================================

\begin{enumerate}
  \item \textbf{Build}: CMake build with \texttt{-DERF\_ENABLE\_MPI=ON} from an
        out-of-source build directory --- should complete without new warnings.
  \item \textbf{Functional}: Add an instrument block to any moist test input file
        (e.g., \texttt{Exec/RegTests/Bomex/inputs\_Bomex}), run 100 steps,
        verify output files exist with monotonically increasing time values.
  \item \textbf{Physics sanity}:
  \begin{itemize}
    \item $\beta_\mathrm{att}$ peaks in the cloud layer and attenuates above it.
    \item Ceilometer cloud base matches visual inspection of $q_c(z)$ profiles.
    \item Doppler lidar $w$ shows sign-changing turbulence in the boundary layer.
    \item LWP from MWR $\approx \sum_k (\rho q_c)_k \Delta z_k$ (computed independently).
    \item MWR $T(z)$ at the surface matches soundings and initial conditions.
  \end{itemize}
  \item \textbf{Microphysics sweep}: Run the same thermodynamic sounding with
        (a)~a bulk scheme (e.g., Kessler or Morrison) and (b)~SBM or SDM;
        compare $\beta_\mathrm{att}$ and LWP time series.
        Bulk results should be smooth functions of LWC; SBM/SDM results will differ
        in cloud-base onset timing and magnitude but should integrate to similar LWP.
\end{enumerate}

% ============================================================
\begin{thebibliography}{99}
\bibitem{Targ1991}
  R.~Targ, M.~J.\ Kavaya, R.~M.\ Huffaker, and R.~L.\ Bowles,
  ``Coherent lidar airborne windshear sensor: performance evaluation,''
  \emph{Appl.\ Opt.}, 30(15), 2013--2026, 1991.
  \doi{10.1364/AO.30.002013}
\bibitem{BohrenHuffman1983}
  C.~F.\ Bohren and D.~R.\ Huffman,
  \emph{Absorption and Scattering of Light by Small Particles},
  Wiley-VCH, 1983.
\bibitem{Rosenkranz1998}
  P.~W.\ Rosenkranz,
  ``Water vapor microwave continuum absorption: A comparison of measurements and models,''
  \emph{Radio Sci.}, 33, 919--928, 1998.
\bibitem{Liebe1993}
  H.~J.\ Liebe, G.~A.\ Hufford, and M.~G.\ Cotton,
  ``Propagation modeling of moist air and suspended water/ice particles at frequencies below 1000 GHz,''
  \emph{Proc. AGARD}, 52nd Specialists' Meeting, 1993.
\bibitem{Ellison2007}
  W.~J.\ Ellison,
  ``Permittivity of Pure Water, at Standard Atmospheric Pressure, over the Frequency Range
  0--25 THz and the Temperature Range 0--100\,$^\circ$C,''
  \emph{J.\ Phys.\ Chem.\ Ref.\ Data}, 36(1), 1--18, 2007.
\end{thebibliography}

\end{document}
